% answers.tex —— 1958 年·四川卷 · 参考答案与详细解析片段
% 由 tools/generate_projects.py 依来源总账自 1958/1958_四川_数学真题_试题.md + 1958/1958_四川_数学真题_解析正文.md + 1958/1958_四川_数学真题_答案要点.md 生成，勿手改（改源 MD 后重跑）


\begin{gktip}
本卷解析由本库不参照答案独立推导（原卷仅题面，无答案页）。两点预告：第 7 题原卷两问以 (I)、(II) 编号，解析按转录体例作（1）（2）；第 5 题按题设推得各边含四次根号，系题设数据的自然结果，非推算错误。
\end{gktip}

\begin{answer}
\gksolhead{1}{}

\textbf{题目}：

化简 $2^{-\frac12}+\cos 225^\circ+\lg 5-\lg\sin 30^\circ+8^{\frac23}$。

\Needspace{4\baselineskip}
\textbf{解析}

\Needspace{5\baselineskip}
逐项化简：

\nopagebreak[4]
\[
2^{-\frac12}=\frac{1}{\sqrt2}=\frac{\sqrt2}{2},\qquad
\cos 225^\circ=\cos(180^\circ+45^\circ)=-\cos 45^\circ=-\frac{\sqrt2}{2}
\]

\nopagebreak[4]
\[
8^{\frac23}=\left(\sqrt[3]{8}\right)^2=2^2=4
\]

\Needspace{5\baselineskip}
对数两项合并：

\nopagebreak[4]
\[
\lg 5-\lg\sin 30^\circ=\lg\frac{5}{\frac12}=\lg 10=1
\]

\Needspace{5\baselineskip}
于是原式

\nopagebreak[4]
\[
2^{-\frac12}+\cos 225^\circ+\lg 5-\lg\sin 30^\circ+8^{\frac23}
=\frac{\sqrt2}{2}-\frac{\sqrt2}{2}+1+4=5
\]

\gkanswer{答案：$5$}

\end{answer}
\begin{answer}
\gksolhead{2}{}

\textbf{题目}：

解不等式 $\dfrac{x+1}{x-1}>1$。

\Needspace{4\baselineskip}
\textbf{解析}

\Needspace{5\baselineskip}
把右边移到左边并通分：

\nopagebreak[4]
\[
\frac{x+1}{x-1}-1>0
\quad\Longleftrightarrow\quad
\frac{(x+1)-(x-1)}{x-1}>0
\quad\Longleftrightarrow\quad
\frac{2}{x-1}>0
\]

\Needspace{5\baselineskip}
分母不为零（$x\neq 1$），故不等式等价于 $x-1>0$，即

\nopagebreak[4]
\[
x>1
\]

\gkanswer{答案：$x>1$}

\end{answer}
\begin{answer}
\gksolhead{3}{}

\textbf{题目}：

证明 $\tan\left(x+\dfrac\pi4\right)+\tan\left(x-\dfrac\pi4\right)=2\tan 2x$。

\Needspace{4\baselineskip}
\textbf{解析}

\Needspace{5\baselineskip}
设 $t=\tan x$。由两角和与两角差的正切公式：

\nopagebreak[4]
\[
\tan\left(x+\frac\pi4\right)=\frac{t+1}{1-t},\qquad
\tan\left(x-\frac\pi4\right)=\frac{t-1}{1+t}
\]

\Needspace{5\baselineskip}
通分相加（第二个分式通分后分子为 $(t-1)(1-t)=-(t-1)^2$）：

\nopagebreak[4]
\[
\tan\left(x+\frac\pi4\right)+\tan\left(x-\frac\pi4\right)
=\frac{(t+1)^2-(t-1)^2}{1-t^2}
=\frac{(t^2+2t+1)-(t^2-2t+1)}{1-t^2}
=\frac{4t}{1-t^2}
\]

\Needspace{5\baselineskip}
而由二倍角公式

\nopagebreak[4]
\[
2\tan 2x=2\cdot\frac{2t}{1-t^2}=\frac{4t}{1-t^2}
\]

两端相等，恒等式成立。以上变形在各表达式有意义的前提下进行，即 $t\neq\pm 1$ 且 $\tan 2x$ 有意义，合起来即 $x\neq\dfrac\pi4+\dfrac{k\pi}{2}$（$k$ 为整数）。

\gkanswer{答案：见解析版}

\end{answer}
\begin{answer}
\gksolhead{4}{}

\textbf{题目}：

若 $\tan\alpha$ 和 $\tan\beta$ 是方程 $x^2+px+q=0$ 的根，用 $p,q$ 来表示 $\tan 2(\alpha+\beta)$。

\Needspace{4\baselineskip}
\textbf{解析}

\Needspace{5\baselineskip}
由韦达定理：

\nopagebreak[4]
\[
\tan\alpha+\tan\beta=-p,\qquad \tan\alpha\,\tan\beta=q
\]

\Needspace{5\baselineskip}
由两角和的正切公式：

\nopagebreak[4]
\[
\tan(\alpha+\beta)=\frac{\tan\alpha+\tan\beta}{1-\tan\alpha\tan\beta}
=\frac{-p}{1-q}
\]

\Needspace{5\baselineskip}
记 $u=\tan(\alpha+\beta)=\dfrac{-p}{1-q}$，由二倍角公式：

\nopagebreak[4]
\[
\tan 2(\alpha+\beta)=\frac{2u}{1-u^2}
=\frac{2\cdot\dfrac{-p}{1-q}}{1-\dfrac{p^2}{(1-q)^2}}
=\frac{-2p(1-q)}{(1-q)^2-p^2}
=\frac{2p(1-q)}{p^2-(1-q)^2}
\]

（表达式有意义的前提：方程有实根要求 $p^2\ge 4q$；上式分母不为零要求 $q\neq 1$ 且 $(1-q)^2\neq p^2$。）

\gkanswer{答案：$\tan 2(\alpha+\beta)=\dfrac{2p(1-q)}{p^2-(1-q)^2}$}

\end{answer}
\begin{answer}
\gksolhead{5}{}

\textbf{题目}：

等腰梯形的上底、腰、下底成等差数列，且底角为 $60^\circ$，梯形的面积为 $2\sqrt2$，求梯形的各边。

\Needspace{4\baselineskip}
\textbf{解析}

\Needspace{5\baselineskip}
设上底为 $a$、腰为 $c$、下底为 $b$。由「上底、腰、下底成等差数列」：

\nopagebreak[4]
\[
2c=a+b
\]

\Needspace{5\baselineskip}
等腰梯形两腰对称，每条腰在下底上的射影长为 $c\cos 60^\circ=\dfrac c2$，故下底

\nopagebreak[4]
\[
b=a+2\cdot\frac c2=a+c
\]

\Needspace{5\baselineskip}
代入 $2c=a+b=2a+c$，得 $c=2a$，从而 $b=a+c=3a$。梯形的高为

\nopagebreak[4]
\[
h=c\sin 60^\circ=2a\cdot\frac{\sqrt3}{2}=\sqrt3\,a
\]

\Needspace{5\baselineskip}
面积为

\nopagebreak[4]
\[
\frac{a+b}{2}\,h=\frac{a+3a}{2}\cdot\sqrt3\,a=2\sqrt3\,a^2
\]

\Needspace{5\baselineskip}
由 $2\sqrt3\,a^2=2\sqrt2$ 得

\nopagebreak[4]
\[
a^2=\frac{\sqrt2}{\sqrt3}=\sqrt{\frac23},
\qquad
a=\sqrt[4]{\frac23}
\]

\Needspace{5\baselineskip}
所以上底、腰、下底分别为

\nopagebreak[4]
\[
\sqrt[4]{\frac23},\qquad
2\sqrt[4]{\frac23},\qquad
3\sqrt[4]{\frac23}
\]

（校验：面积 $=2\sqrt3\left(\sqrt[4]{\dfrac23}\right)^2=2\sqrt3\sqrt{\dfrac23}=2\sqrt2$，与题设相符。）

\gkanswer{答案：上底 $\sqrt[4]{\dfrac23}$，腰 $2\sqrt[4]{\dfrac23}$，下底 $3\sqrt[4]{\dfrac23}$}

\end{answer}
\begin{answer}
\gksolhead{6}{}

\textbf{题目}：

解方程 $9^x-10\cdot 3^{x+1}+81=0$。

\Needspace{4\baselineskip}
\textbf{解析}

\Needspace{5\baselineskip}
令 $t=3^x$（$t>0$），则 $9^x=t^2$，$3^{x+1}=3t$，方程化为

\nopagebreak[4]
\[
t^2-30t+81=0
\]

\Needspace{5\baselineskip}
判别式 $\Delta=(-30)^2-4\cdot 81=900-324=576=24^2$，故

\nopagebreak[4]
\[
t=\frac{30\pm 24}{2}=27\quad\text{或}\quad 3
\]

由 $3^x=27$ 得 $x=3$；由 $3^x=3$ 得 $x=1$。

\gkanswer{答案：$x=3$ 或 $x=1$}

\end{answer}
\begin{answer}
\gksolhead{7}{}

\textbf{题目}：

同底同侧的两个三角形 $ABC$ 和 $ABD$，$AC$、$BC$、$BD$、$AD$ 的中点分别为 $M$、$N$、$P$、$Q$，求证：

（1） $MNPQ$ 为一平行四边形；

（2） 平行四边形 $MNPQ$ 的面积等于两三角形面积差的一半。

\begin{gktip}
\textbf{题号说明}：原卷本题两问以 (I)、(II) 编号，本转录按体例统一为（1）、（2），内容不变。
\end{gktip}

\Needspace{4\baselineskip}
\textbf{解析}

两个三角形公共底为 $AB$，$C$、$D$ 在 $AB$ 的同侧；$M$、$N$、$P$、$Q$ 分别是 $AC$、$BC$、$BD$、$AD$ 的中点。

\gksubpart{（1）}

\Needspace{5\baselineskip}
在 $\triangle ABC$ 中，$M$、$N$ 分别是 $AC$、$BC$ 的中点，由中位线定理：

\nopagebreak[4]
\[
MN\parallel AB,\qquad MN=\frac12 AB
\]

\Needspace{5\baselineskip}
在 $\triangle ABD$ 中，$Q$、$P$ 分别是 $AD$、$BD$ 的中点，同理：

\nopagebreak[4]
\[
QP\parallel AB,\qquad QP=\frac12 AB
\]

于是 $MN\parallel QP$ 且 $MN=QP$。一组对边平行且相等的四边形是平行四边形，故 $MNPQ$ 为平行四边形。

\gksubpart{（2）}

\Needspace{5\baselineskip}
设 $C$、$D$ 到直线 $AB$ 的距离分别为 $h_C$、$h_D$。$MN$ 是 $\triangle ABC$ 的中位线，所在直线平行于 $AB$ 且到 $AB$ 的距离为 $\dfrac{h_C}{2}$；同理 $QP$ 所在直线到 $AB$ 的距离为 $\dfrac{h_D}{2}$。于是平行四边形 $MNPQ$ 以 $MN=\dfrac{AB}{2}$ 为底，以两条平行线间的距离 $\left|\dfrac{h_C}{2}-\dfrac{h_D}{2}\right|=\dfrac{|h_C-h_D|}{2}$ 为高，其面积为

\nopagebreak[4]
\[
S_{MNPQ}=\frac{AB}{2}\cdot\frac{|h_C-h_D|}{2}=\frac{AB\,|h_C-h_D|}{4}
\]

\Needspace{5\baselineskip}
而同底同侧的两三角形面积差为

\nopagebreak[4]
\[
\left|S_{\triangle ABC}-S_{\triangle ABD}\right|
=\left|\frac{AB\cdot h_C}{2}-\frac{AB\cdot h_D}{2}\right|
=\frac{AB\,|h_C-h_D|}{2}
\]

恰为 $S_{MNPQ}$ 的 $2$ 倍。故平行四边形 $MNPQ$ 的面积等于两三角形面积差的一半。

\gkanswer{答案（1）：见解析版}

\gkanswer{答案（2）：见解析版}

\end{answer}
\begin{answer}
\gksolhead{8}{}

\textbf{题目}：

某圆锥形的人造卫星中，有一个内切球形仓，已知此球的体积为圆锥体积的一半，求圆锥母线与底面的夹角。

\Needspace{4\baselineskip}
\textbf{解析}

\Needspace{5\baselineskip}
设圆锥底面半径为 $R$、高为 $H$、母线为 $l$，内切球半径为 $r$。轴截面是底为 $2R$、高为 $H$、腰为 $l$ 的等腰三角形，内切圆半径即球半径，由「面积 $=$ 内切圆半径 $\times$ 半周长」：

\nopagebreak[4]
\[
r=\frac{\frac12\cdot 2R\cdot H}{R+l}=\frac{RH}{R+l}
\]

\Needspace{5\baselineskip}
球体积为圆锥体积的一半：

\nopagebreak[4]
\[
\frac43\pi r^3=\frac12\cdot\frac13\pi R^2H
\quad\Longleftrightarrow\quad
8r^3=R^2H
\]

\Needspace{5\baselineskip}
代入 $r=\dfrac{RH}{R+l}$：

\nopagebreak[4]
\[
8\cdot\frac{R^3H^3}{(R+l)^3}=R^2H
\quad\Longleftrightarrow\quad
8RH^2=(R+l)^3
\]

\Needspace{5\baselineskip}
设母线与底面的夹角为 $\theta$，则 $R=l\cos\theta$，$H=l\sin\theta$，代入上式并约去 $l^3$：

\nopagebreak[4]
\[
8\sin^2\theta\cos\theta=(1+\cos\theta)^3
\]

\Needspace{5\baselineskip}
记 $c=\cos\theta$（$0<c<1$），由 $1-c^2=(1-c)(1+c)$：

\nopagebreak[4]
\[
8c(1-c)(1+c)=(1+c)^3
\]

\Needspace{5\baselineskip}
两边除以正数 $(1+c)^2$：

\nopagebreak[4]
\[
8c(1-c)=(1+c)^2
\]

\Needspace{5\baselineskip}
展开整理：

\nopagebreak[4]
\[
8c-8c^2=1+2c+c^2
\quad\Longleftrightarrow\quad
9c^2-6c+1=0
\quad\Longleftrightarrow\quad
(3c-1)^2=0
\]

\Needspace{5\baselineskip}
故 $c=\dfrac13$，即

\nopagebreak[4]
\[
\cos\theta=\frac13,
\qquad
\theta=\arccos\frac13
\]

（方程在 $0<c<1$ 内仅此一根（二重根），夹角唯一。）

\gkanswer{答案：母线与底面的夹角 $\theta=\arccos\dfrac13$}

\end{answer}